\documentclass[11pt,a4paper]{article}
\usepackage[margin=22mm]{geometry}
\usepackage{fontspec,babel}
\babelprovide[main,import]{russian}
\babelfont{rm}{Times New Roman}
\babelfont{sf}{Arial}
\usepackage{amsmath,booktabs,graphicx,xcolor}
\usepackage[colorlinks=true,linkcolor=teal,urlcolor=teal]{hyperref}
\setlength{\parindent}{0pt}
\setlength{\parskip}{6pt}
\newcommand{\heading}[1]{\section*{\sffamily #1}}
\newcommand{\plot}[1]{\begin{center}\includegraphics[width=\linewidth]{#1}\end{center}}
\newcommand{\repo}{https://github.com/snejokzboxik/three-phase-12-electrode-trap}
\begin{document}

{\sffamily\LARGE\bfseries Эффективная динамика и удержание\\[3pt]
в трёхфазной 12-электродной ловушке\par}
\vspace{4mm}
{\sffamily Численное исследование идеальной геометрии \hfill 8 ноября 2026 г.}
\vspace{3mm}\hrule

\heading{1. Реконструкция медленной динамики}
На синтетических траекториях нескольких амплитуд восстановлен
аналитический закон \(B_6r^6\). Ошибка показателя составляет
\(1{,}2\times10^{-13}\), относительная ошибка коэффициента
\(7{,}1\times10^{-10}\). Стандартный FEM-псевдопотенциал B
воспроизводит радиальное движение, а модель E дополнительно учитывает
перпендикулярное скорости ускорение \(\gamma(r)J\dot{\mathbf R}\), где
\[
J(x,y)=(-y,x),\qquad
\gamma(r)=\sigma\,1{,}26760747\,\mathrm{s}^{-1}(r/1\,\mathrm{mm})^4.
\]
Знак \(\sigma\) меняется при перестановке фаз ABC и ACB.
Коэффициент был зафиксирован до независимого расчёта прямых P4-траекторий.

\begin{center}
\small
\begin{tabular}{ccrrrr}
\toprule
\(r_0\), мм & Фазы & P4 h35, град/с & E, град/с & RMS B, \% & RMS E, \%\\
\midrule
0,75 & ABC & \(+3{,}53248\) & \(+3{,}53331\) & 0,234146 & 0,000125\\
0,75 & ACB & \(-3{,}53694\) & \(-3{,}53331\) & 0,234443 & 0,000258\\
1,25 & ABC & \(+27{,}24813\) & \(+27{,}26324\) & 1,083732 & 0,001623\\
1,25 & ACB & \(-27{,}32063\) & \(-27{,}26324\) & 1,086634 & 0,002538\\
\bottomrule
\end{tabular}
\end{center}

Векторная RMS нормирована на 1 мм, горизонт составляет 3,25 медленного
периода. Межсеточное отличие траекторий h35 и h50 составляет
0,023--0,038\% и задаёт отдельный масштаб чувствительности. Глобальная
сходимость прямого расчёта за 50 медленных периодов остаётся открытой.

\plot{outputs/effective_reconstruction/p4_transfer/v1/report/p4_transfer.png}
\smallskip
{\small Рис. 1. Прецессия, ошибки и поперечное движение на независимых P4-стартах.}

\clearpage
\heading{2. Удержание в зависимости от стартового радиуса}

Карта при \(a=0\) сравнивает N48, гармоническое представление поля P4,
с идеальным квадруполем на 604 общих узлах. В каждом узле проверено по
шесть стартов с общим горизонтом 10\,000 RF-периодов и апертурой 8 мм.
Все шесть удержались в 248 узлах N48 и 91 узле квадруполя. За первой
границей Матьё \( |q_{\rm M}^{\rm ref}|\approx0{,}908 \) таких узлов
у N48 было 91, а у квадруполя ушли все проверенные старты.

\plot{outputs/stability_p4/matched_radius_v2/radius_points.png}
{\small Рис. 2. Рассчитанные узлы. Нижние панели увеличивают окрестность
первой границы Матьё.}

При опорном \( |q_{\rm M}^{\rm ref}|\approx0{,}03958 \) переход от полного
к частичному удержанию заключён между 6,75 и 6,8125 мм для N48 и между
4,5 и 5,0 мм для квадруполя. Границы взяты по соседним расчётным узлам.
Число узлов зависит от плотности адаптивной сетки.

\clearpage
\heading{3. Интерполяция радиальной карты}
\plot{outputs/stability_p4/matched_radius_v2/radius_contours.png}
{\small Рис. 3. Контуры доли удержанных стартов. Для интерполяции
использована триангуляция в \((\log_{10}|q_{\rm M}^{\rm ref}|,r_0)\).
Разреженные участки оставлены белыми. На панели квадруполя данные по
разные стороны первой границы Матьё интерполированы отдельно.}

В двух стартах квадруполя изменение допусков сдвинуло время
зарегистрированного выхода примерно на 0,51 RF-периода. Классификация
всех 3624 стартов сохранилась. Сравнение по радиусу характеризует
поперечное удержание на конечном времени.

\clearpage
\heading{4. Сравнение в общей плоскости \((a,q_{\rm M}^{\rm ref})\)}

В обеих моделях наложен идеальный квадрупольный DC-потенциал
\(\Phi_{\rm DC}=U_{\rm DC}(x^2-y^2)/R^2\), \(R=8\) мм. Тогда
\[
a=\frac{8QU_{\rm DC}}{mR^2\Omega^2},\qquad
|q_{\rm M}^{\rm ref}|=\frac{4|Q|U_0}{mR^2\Omega^2}.
\]
Рассчитано 120 узлов с шестью стартами на радиусе 1 мм на протяжении
1000 RF-периодов. При \(|q_{\rm M}^{\rm ref}|=0{,}706\) N48 теряет все
старты при \(a=0{,}025\), а квадруполь удерживает все при \(a=0{,}160\).
При \(|q_{\rm M}^{\rm ref}|=0{,}9973\) и \(a=0{,}010\) N48 удерживает
6 из 6, квадруполь 0 из 6.

\plot{outputs/stability_p4/aq_ideal_dc_v1/RF1000/aq_contours_article.png}
{\small Рис. 4. Интерполированные уровни 50\% и 5 из 6. Крест отмечает
точку, где два исхода N48 изменились при проверке шага. Окрестность
точки исключена из интерполяции.}

Вблизи центра у N48 ведущая RF-сила порядка \(r^5\), а выталкивающая
DC-сила линейна. У квадруполя обе силы линейны в приближении
псевдопотенциала. Это объясняет различную чувствительность к DC около
центра. Границы карты зависят от ансамбля и горизонта. DC-поле реальных
электродов в этой серии заменено идеальным.

\clearpage
\heading{5. Локальная устойчивость и применение}

Для 11 значений связи найдены 66 замкнутых орбит на шести ветвях N48.
Все найденные ветви линейно неустойчивы. Для трёх орбит этот класс
подтверждён на прямых полях P4 h35 и h50. Наблюдаемая локальная
неустойчивость сосуществует с конечновременным удержанием частицы.
Поиск охватывает выбранные семейства периодических решений.

\plot{outputs/floquet_orbits/v1/floquet_branches.png}
{\small Рис. 5. Максимальный модуль мультипликатора Флоке найденных ветвей.}

\textbf{Связь с масс-спектрометрией.} При фиксированных напряжениях и
частоте частицы одного знака заряда с разным \(m/z\) располагаются
вдоль одной прямой в плоскости \((a,q_{\rm M}^{\rm ref})\). Широкое
поперечное удержание при слабом DC мотивирует проверку трёхфазной
геометрии как ионовода. Для вывода о пропускании или разрешении нужна
кривая \(T(m/z)\) в приборе конечной длины с одинаковым ансамблем
ввода и краевыми полями. Обычный двухфазный октуполь нужен как эталон
для выделения эффекта трёхфазного поля.

\textbf{Данные.} Таблицы CSV, конфигурации, графики и хеши размещены
в \href{\repo}{репозитории проекта}. Тяжёлые массивы и рабочие
контрольные точки хранятся локально. Литературный прецедент для
октупольных карт \href{https://doi.org/10.1016/0168-1176(86)80003-9}{Hägg и Szabó, 1986}.

\end{document}
